\documentclass[12pt,fullpage]{article}

\usepackage{amsthm}
\usepackage{amssymb}
\usepackage[fleqn]{amsmath}
\usepackage{enumitem}

\newtheorem{theorem}{Theorem}[section]
\newtheorem{lemma}[theorem]{Lemma}
\newtheorem{proposition}[theorem]{Proposition}
\newtheorem{corollary}[theorem]{Corollary}

\theoremstyle{definition}
\newtheorem{definition}[theorem]{Definition}
\newtheorem{remark}[theorem]{Remark}

\newcommand{\A}{\mathcal{A}}
\newcommand{\B}{\mathcal{B}}
\newcommand{\N}{\mathbb{N}}
\newcommand{\Z}{{\mathbb{Z}}}
\newcommand{\AUT}{\operatorname{AUT}}
\renewcommand{\phi}{\varphi}

\begin{document}

\section*{Sub-lemma (G2): $\AUT(\A)$ is countable --- in fact $\AUT(\A)\cong\Z$}

This section discharges the framing hypothesis of Theorems~2.2 and~2.3 by
giving an \emph{explicit} structure $\A$, computing its automorphism group
completely, and showing it is countable.  We then verify that the very same
$\A$ has the structural features (K1)--(K4) used in the rest of the argument
(in particular by Sub-lemma~G5 below), so that establishing countability does
\emph{not} rigidify $\A$ past the point where the diagonal/negative clause
remains available.

The critic's worry is the correct one to address: a naive family of infinitely
many \emph{independent} binary (orientation) choices produces
$(\Z/2)^{\N}$, which is uncountable.  The construction below installs a
\emph{rigid coordinating spine} that forces every automorphism to act by the
same translation amount on every gadget and forbids all reflections globally,
collapsing what would have been $(\Z/2)^{\N}\rtimes(\ \cdots\ )$ down to the
single copy of $\Z$ of simultaneous translations.  The spine is, crucially,
weak enough that a \emph{single} unanchored gadget still admits a local
orientation reversal (property (K3)), which is exactly what the negative
clause needs.

\subsection*{The construction}

\begin{definition}[The structure $\A$]\label{def:A}
The universe is
\[
A \;=\; \{\,(n,i) : n\in\N,\ i\in\Z\,\},
\]
a set in explicit bijection with $\N$ (e.g.\ via a computable pairing of
$\N\times\Z$ with $\N$); thus $\A$ is computably represented.  For each $n$ put
\[
G_n \;=\; \{\,(n,i) : i\in\Z\,\},
\]
so $A=\bigsqcup_{n\in\N}G_n$ is partitioned into the \emph{gadgets} $G_n$.
The language has three binary relation symbols, interpreted as follows.
\begin{itemize}
\item \textbf{Within-gadget edges.}
  \[
  P\bigl((n,i),(m,j)\bigr)\ \Longleftrightarrow\ n=m\ \text{and}\ |i-j|=1 .
  \]
  Thus on each gadget $G_n$, the relation $P$ is the edge relation of an
  \emph{undirected} bi-infinite path.
\item \textbf{Anchor orientation on $G_0$.}
  \[
  D\bigl((n,i),(m,j)\bigr)\ \Longleftrightarrow\ n=m=0\ \text{and}\ j=i+1 .
  \]
  Thus $D$ holds only inside the single gadget $G_0$, where it is the edge
  relation of a \emph{directed} bi-infinite path (a $\Z$-line with successor).
  On every gadget $G_n$ with $n\ge 1$, $D$ is empty.
\item \textbf{Coordinating spine (rungs).}
  \[
  R\bigl((n,i),(m,j)\bigr)\ \Longleftrightarrow\ m=n+1\ \text{and}\ j=i .
  \]
  Thus $R$ is a \emph{directed} relation joining the point $(n,i)$ of gadget
  $G_n$ to the point $(n+1,i)$ of gadget $G_{n+1}$ at the \emph{same} position
  $i$.  The reduct $(A;R)$ is, for each fixed $i$, a directed ray
  $(0,i)\to(1,i)\to(2,i)\to\cdots$, and these rays are vertex-disjoint.
\end{itemize}
\end{definition}

Three design features carry the whole argument, and we name them now:
the \emph{rungs} $R$ couple all gadgets (forcing a uniform translation), their
\emph{direction} together with the indexing $n\in\N$ pins each gadget index
(no gadget is permuted), and the \emph{directed anchor} $D$ on $G_0$ destroys
the global reflection.

\subsection*{Computation of $\AUT(\A)$}

For $k\in\Z$ define $\sigma_k\colon A\to A$ by
\[
\sigma_k(n,i) \;=\; (n,\,i+k).
\]

\begin{lemma}[Each $\sigma_k$ is an automorphism]\label{lem:shifts-are-auts}
For every $k\in\Z$, $\sigma_k\in\AUT(\A)$, and $\sigma_k\circ\sigma_\ell=
\sigma_{k+\ell}$.  Hence $k\mapsto\sigma_k$ is a homomorphism $\Z\to\AUT(\A)$.
\end{lemma}

\begin{proof}
$\sigma_k$ is a bijection of $A$ with inverse $\sigma_{-k}$, and
$\sigma_k\sigma_\ell=\sigma_{k+\ell}$ is immediate from
$(i+\ell)+k=i+(k+\ell)$.  We check each relation is preserved (and, applying
the same computation to $\sigma_{-k}$, reflected).

\emph{$P$.} $P((n,i),(m,j))$ iff $n=m$ and $|i-j|=1$; then
$\sigma_k(n,i)=(n,i+k)$, $\sigma_k(m,j)=(m,j+k)$, and indeed $n=m$ and
$|(i+k)-(j+k)|=|i-j|=1$, so $P(\sigma_k(n,i),\sigma_k(m,j))$ holds; conversely
the same with $-k$.

\emph{$D$.} $D((n,i),(m,j))$ iff $n=m=0$ and $j=i+1$; then the images are
$(0,i+k),(0,j+k)$ with $j+k=(i+k)+1$, so $D$ is preserved; conversely with
$-k$.

\emph{$R$.} $R((n,i),(m,j))$ iff $m=n+1$ and $j=i$; the images are
$(n,i+k),(n+1,i+k)$, still with second index $n+1=n+1$ and equal positions
$i+k=i+k$, so $R$ is preserved; conversely with $-k$.

Thus $\sigma_k$ preserves and reflects all three relations, i.e.\
$\sigma_k\in\AUT(\A)$.
\end{proof}

\begin{lemma}[Every automorphism fixes each gadget index]\label{lem:fix-index}
Let $\phi\in\AUT(\A)$.  Then for every $(n,i)\in A$ there is $i'\in\Z$ with
$\phi(n,i)=(n,i')$.  Equivalently, every gadget $G_n$ is $\phi$-invariant
(setwise), and $\phi$ preserves the gadget index $n$.
\end{lemma}

\begin{proof}
Write $\phi(n,i)=(N(n,i),\,I(n,i))$.  We first identify the gadget index from
the relation $R$ alone.  Consider the binary relation $R$ as a directed graph
on $A$.  For a point $p\in A$ define its \emph{in-degree} and
\emph{out-degree} in $R$.  From Definition~\ref{def:A}:
\[
\text{out-degree}_R(n,i)=1\ \text{(the edge to }(n+1,i)),\qquad
\text{in-degree}_R(n,i)=
\begin{cases}
0,& n=0,\\[2pt]
1,& n\ge 1\ \text{(the edge from }(n-1,i)).
\end{cases}
\]
Automorphisms preserve in- and out-degrees in each relation.  Hence:

\smallskip
\noindent\emph{Step 1: $\phi$ preserves $G_0$ setwise.}
$(n,i)\in G_0$ iff its $R$-in-degree is $0$.  Since $\phi$ preserves
$R$-in-degree, $\phi(G_0)=G_0$.

\smallskip
\noindent\emph{Step 2: $\phi$ preserves each $G_n$ setwise, by induction on
$n$.}  Suppose $\phi(G_m)=G_m$ for all $m\le n$; we show $\phi(G_{n+1})=
G_{n+1}$.  A point $q$ lies in $G_{n+1}$ iff there is a (necessarily unique)
point $p$ with $R(p,q)$ and $p\in G_n$.  Indeed $R(p,q)$ holds iff
$q=(n+1,i)$ and $p=(n,i)$ for some $i$; so $q\in G_{n+1}$ iff its unique
$R$-predecessor lies in $G_n$.  Take $q\in G_{n+1}$, with $R$-predecessor
$p\in G_n$.  Then $R(\phi(p),\phi(q))$ holds (as $\phi$ preserves $R$), and
$\phi(p)\in G_n$ by the induction hypothesis.  The unique $R$-successor of any
point of $G_n$ lies in $G_{n+1}$; hence $\phi(q)\in G_{n+1}$.  This shows
$\phi(G_{n+1})\subseteq G_{n+1}$, and applying the same argument to
$\phi^{-1}$ gives equality.  By induction $\phi(G_n)=G_n$ for all $n$.

Thus $\phi$ preserves the gadget index, proving the claim.
\end{proof}

\begin{lemma}[The position shift is uniform across gadgets]\label{lem:uniform}
Let $\phi\in\AUT(\A)$.  There is a single bijection $f\colon\Z\to\Z$ such that
\[
\phi(n,i)=(n,\,f(i))\qquad\text{for all }n\in\N,\ i\in\Z .
\]
\end{lemma}

\begin{proof}
By Lemma~\ref{lem:fix-index}, $\phi(n,i)=(n,I(n,i))$ for some
$I(n,i)\in\Z$.  Fix $i\in\Z$ and consider the $R$-ray
$(0,i)\to(1,i)\to(2,i)\to\cdots$.  Since $R(p,q)\Leftrightarrow
R(\phi(p),\phi(q))$ and $\phi$ fixes gadget indices, applying $\phi$ to the
edge $R((n,i),(n+1,i))$ gives $R\bigl((n,I(n,i)),(n+1,I(n+1,i))\bigr)$, which
by the definition of $R$ forces $I(n+1,i)=I(n,i)$.  Hence $I(n,i)$ is
independent of $n$; call it $f(i):=I(0,i)$.  Then
$\phi(n,i)=(n,f(i))$ for all $n,i$.  Finally $f$ is a bijection of $\Z$
because $\phi$ is a bijection of $A$ fixing each gadget index (its restriction
to $G_0\cong\Z$ is exactly $i\mapsto f(i)$).
\end{proof}

\begin{lemma}[The uniform shift must be a translation]\label{lem:translation}
The bijection $f$ of Lemma~\ref{lem:uniform} satisfies $f(i)=i+k$ for a fixed
$k\in\Z$.  Consequently $\phi=\sigma_k$.
\end{lemma}

\begin{proof}
\emph{$f$ preserves the undirected path on $\Z$.}  Restrict $\phi$ to $G_1$
(where $P$ is the undirected bi-infinite path and $D$ is empty).  Since $\phi$
preserves $P$ and acts on $G_1$ by $i\mapsto f(i)$, we have for all $i$:
\[
|i-(i+1)|=1\ \Rightarrow\ |f(i)-f(i+1)|=1 .
\]
So $f(i+1)\in\{f(i)+1,f(i)-1\}$ for every $i$.  A bijection $f\colon\Z\to\Z$
with $|f(i+1)-f(i)|=1$ for all $i$ is a graph automorphism of the bi-infinite
path; such an $f$ is monotone (if it changed direction at some $i$, i.e.\
$f(i+1)=f(i)+1$ and $f(i+2)=f(i+1)-1=f(i)$, then $f(i)=f(i+2)$ contradicting
injectivity, and symmetrically for the other turn), hence either
$f(i)=i+k$ (orientation-preserving) or $f(i)=c-i$ (orientation-reversing)
for fixed $k$ or $c$ in $\Z$.

\emph{The reversal is excluded by the anchor $D$ on $G_0$.}  Suppose
$f(i)=c-i$ for some $c\in\Z$.  On $G_0$ the relation $D$ is the directed
successor: $D((0,i),(0,i+1))$ holds for all $i$.  Since $\phi$ preserves $D$
and acts on $G_0$ by $i\mapsto f(i)$, applying $\phi$ to this edge yields
\[
D\bigl((0,f(i)),(0,f(i+1))\bigr),\qquad\text{i.e.}\qquad f(i+1)=f(i)+1 .
\]
But $f(i)=c-i$ gives $f(i+1)=c-i-1=f(i)-1\ne f(i)+1$, a contradiction.
Hence $f$ is orientation-preserving, $f(i)=i+k$ for a fixed $k\in\Z$, and
therefore $\phi=\sigma_k$ by Lemma~\ref{lem:uniform}.
\end{proof}

\begin{theorem}[Complete characterization]\label{thm:aut}
The map $k\mapsto\sigma_k$ is a group isomorphism
\[
\Z\ \xrightarrow{\ \sim\ }\ \AUT(\A),\qquad k\mapsto\sigma_k .
\]
In particular $\AUT(\A)\cong\Z$ is infinite cyclic, hence \emph{countable},
generated by $\sigma:=\sigma_1$, and the structure $\A$ has nontrivial
automorphisms.
\end{theorem}

\begin{proof}
By Lemma~\ref{lem:shifts-are-auts}, $k\mapsto\sigma_k$ is a homomorphism
$\Z\to\AUT(\A)$; it is injective because $\sigma_k=\mathrm{id}$ forces
$i+k=i$, i.e.\ $k=0$.  By Lemmas~\ref{lem:fix-index}, \ref{lem:uniform},
and~\ref{lem:translation}, every $\phi\in\AUT(\A)$ equals some $\sigma_k$, so
the homomorphism is surjective.  Hence it is an isomorphism, $\AUT(\A)\cong\Z$,
which is countably infinite.
\end{proof}

\begin{remark}[How the uncountable trap is avoided]\label{rem:trap}
Without the rungs $R$, the gadgets would be independent and each would
contribute its own $D_\infty$ (translations and a reflection) acting
independently, giving an automorphism group containing
$(\Z/2)^{\N}$ (the independent reflections) --- uncountable.  The rungs are a
\emph{rigid, oriented} coordinating spine: Lemma~\ref{lem:fix-index} uses
their orientation and the well-ordered index $n\in\N$ (the source $G_0$ has
$R$-in-degree $0$) to fix every gadget setwise, and Lemma~\ref{lem:uniform}
uses their equal-position joining to force \emph{one} shift $f$ on all gadgets
at once.  The remaining binary choice --- the single reflection still available
to a uniform $f$ --- is killed once and for all by the directed anchor $D$ on
the lone gadget $G_0$ (Lemma~\ref{lem:translation}).  What survives is exactly
the diagonal $\{\sigma_k:k\in\Z\}\cong\Z$ inside $(\Z\rtimes\Z/2)$, i.e.\ the
group of \emph{simultaneous} translations.  This is the ``direct sum, not
direct product'' phenomenon the critic anticipated, realized in its strongest
form: not merely finite support, but a single common $k$.
\end{remark}

\subsection*{Verification of the structural features (K1)--(K4)}

We now confirm that the explicit $\A$ of Definition~\ref{def:A} satisfies the
properties (K1)--(K4) invoked by Sub-lemma~G5, and in particular that
establishing countability has \emph{not} over-rigidified $\A$.

\begin{proposition}[(K1) holds]\label{prop:K1}
$\AUT(\A)\cong\Z=\langle\sigma\rangle$ with $\sigma=\sigma_1$; each gadget
$G_n$ is $\sigma$-invariant; and $\sigma$ acts on $G_n$ as the fixed-point-free
$\Z$-translation $(n,i)\mapsto(n,i+1)$.  Hence every nontrivial $\sigma^k$
($k\ne0$) moves every point of every gadget, and $\A$ satisfies the hypotheses
of Theorem~2.2 (countable automorphism group).
\end{proposition}

\begin{proof}
Immediate from Theorem~\ref{thm:aut} and the formula
$\sigma^k=\sigma_k$, $\sigma_k(n,i)=(n,i+k)$, which has no fixed point on any
$G_n$ when $k\ne 0$.
\end{proof}

\begin{proposition}[(K2) holds]\label{prop:K2}
For each $n$ and each $k\in\Z$ there is an existential formula
$\theta_n^{(k)}(u,x,y)$ such that, for any fixed base point $b\in G_n$,
$\theta_n^{(k)}(b,\cdot,\cdot)$ defines on $G_n$ the graph of
$\sigma^k\restriction G_n$.
\end{proposition}

\begin{proof}
The relation $P$ restricted to $G_n$ is the edge relation of the undirected
path; ``$x$ and $y$ lie in the same gadget at $P$-distance exactly $|k|$ on
the side of $y$ determined by $u$'' is expressible existentially once a base
point orients the path.  Concretely, fix $k\ge 1$ (the case $k\le -1$ is
symmetric, $k=0$ is ``$x=y$'').  Let
\[
\theta_n^{(k)}(u,x,y)\;:=\;\exists z_0\,\exists z_1\cdots\exists z_{k}\;
\Bigl(z_0=x\ \wedge\ z_k=y\ \wedge\ \bigwedge_{t=0}^{k-1}P(z_t,z_{t+1})\ \wedge\ \mathrm{Or}(u,z_0,z_1)\Bigr),
\]
where $\mathrm{Or}(u,z_0,z_1)$ is the existential formula asserting that the
step $z_0\to z_1$ goes in the direction \emph{away from} the base point $u$,
i.e.\ that the $P$-distance from $u$ to $z_1$ exceeds the $P$-distance from
$u$ to $z_0$ (expressible existentially by exhibiting a $P$-path of suitable
length, using that $u,z_0,z_1$ are colinear on the path).  Once $u=b$ is a
fixed point of $G_n$, the orientation $\mathrm{Or}(b,\cdot,\cdot)$ singles out
the same direction as $\sigma$, and $\theta_n^{(k)}(b,x,y)$ holds on $G_n$
exactly when $y=\sigma^k(x)$.  (The single point $b$ supplies precisely the
orientation that the undirected $P$ lacks intrinsically on $G_n$, $n\ge 1$;
on $G_0$ the relation $D$ already provides orientation, so $b$ may even be
omitted there.)
\end{proof}

\begin{proposition}[(K3) holds]\label{prop:K3}
Let $\bar c$ be any finite subset of $A$.  Then there is a gadget index
$n_0\ge 1$ with $G_{n_0-1},G_{n_0},G_{n_0+1}$ all disjoint from $\bar c$, and
for any such $n_0$ the induced substructure on $\bar c\cup G_{n_0}$ admits an
automorphism $\rho$ that fixes $\bar c$ pointwise, fixes every
$G_m$ ($m\ne n_0$) pointwise, and reverses the orientation of $G_{n_0}$:
\[
\rho(n_0,i)=(n_0,\,-i),\qquad \rho\ \text{the identity elsewhere on its domain.}
\]
In particular $\rho\,\sigma\restriction G_{n_0}\,\rho^{-1}=
\sigma^{-1}\restriction G_{n_0}$.
\end{proposition}

\begin{proof}
$\bar c$ meets only finitely many gadgets, so some $n_0\ge 1$ has
$G_{n_0-1},G_{n_0},G_{n_0+1}$ all disjoint from $\bar c$.  Consider the
induced relational structure $\mathcal{S}$ on the domain
$\bar c\cup G_{n_0}$.  Its relation-instances touching $G_{n_0}$ are:
the $P$-edges \emph{within} $G_{n_0}$ (an undirected path, since $n_0\ge1$ so
$D$ is empty on $G_{n_0}$); no $D$-instances (as $n_0\ge 1$); and any
$R$-instances with one endpoint in $G_{n_0}$.  An $R$-edge incident to
$G_{n_0}$ has its other endpoint in $G_{n_0-1}$ or $G_{n_0+1}$; but those
gadgets are disjoint from $\bar c$, hence \emph{absent} from the domain
$\bar c\cup G_{n_0}$.  Therefore $\mathcal{S}$ contains \emph{no} $R$-edge
incident to $G_{n_0}$.  Consequently the only structure on $G_{n_0}$ inside
$\mathcal{S}$ is the undirected path $P\restriction G_{n_0}$, and there is no
relation linking $G_{n_0}$ to $\bar c$ at all.

Define $\rho$ to be the identity on $\bar c$ and the reflection
$(n_0,i)\mapsto(n_0,-i)$ on $G_{n_0}$.  This is a bijection of the domain
$\bar c\cup G_{n_0}$.  It preserves $P\restriction G_{n_0}$ because
$|i-j|=1\Leftrightarrow|(-i)-(-j)|=1$; it preserves the (empty) $D$-relation
trivially; and it preserves the $R$-relation of $\mathcal{S}$ trivially, since
$\mathcal{S}$ has no $R$-edge touching $G_{n_0}$ and $\rho$ fixes everything
outside $G_{n_0}$.  Hence $\rho\in\AUT(\mathcal{S})$, it fixes $\bar c$
pointwise, and on $G_{n_0}$ it satisfies
$\rho(\sigma(n_0,i))=\rho(n_0,i+1)=(n_0,-i-1)=\sigma^{-1}(n_0,-i)=
\sigma^{-1}(\rho(n_0,i))$, i.e.\ $\rho\sigma\rho^{-1}=\sigma^{-1}$ on
$G_{n_0}$, as claimed.
\end{proof}

\begin{remark}[Why countability did not destroy (K3)]\label{rem:K3-survives}
This is the crux of the critic's concern (b).  The directed anchor $D$ lives
\emph{only} on $G_0$; on every gadget $G_n$ with $n\ge 1$ the within-gadget
relation is the \emph{undirected} path, which is reflection-symmetric.  The
spine $R$ is what would, in the full structure, forbid reflecting a single
gadget --- but in the \emph{induced} substructure on $\bar c\cup G_{n_0}$ the
spine edges incident to $G_{n_0}$ vanish (their other endpoints are not in the
domain).  Thus the global rigidity that yields $\AUT(\A)\cong\Z$ coexists with
local flexibility of any individual unanchored gadget $n_0\ge1$.  This is
exactly the resource asymmetry exploited by Sub-lemma~G5.
\end{remark}

\begin{proposition}[(K4) holds]\label{prop:K4}
For each fixed $k\ne 0$ and each existential formula $\psi(\bar c,x,y)$ whose
parameters $\bar c$ are disjoint from a gadget $G_{n_0}$ (with $G_{n_0\pm1}$
also disjoint from $\bar c$), $\psi$ cannot define $\sigma^k\restriction
G_{n_0}$.  Equivalently, the orientation needed to certify ``move by $k$''
inside $G_{n_0}$ is not available from parameters outside $G_{n_0}$.
\end{proposition}

\begin{proof}
If $\psi(\bar c,x_0,\sigma^k x_0)$ held for some $x_0\in G_{n_0}$, then
applying the orientation-reversing automorphism $\rho$ of
Proposition~\ref{prop:K3} (which fixes $\bar c$ and preserves existential
formulas) would give $\psi(\bar c,\rho x_0,\rho\sigma^k x_0)$ with
$\rho\sigma^k\rho^{-1}=\sigma^{-k}$ on $G_{n_0}$; this would force $\psi$ to
also relate $\rho x_0$ to $\sigma^{-k}(\rho x_0)\ne\sigma^{k}(\rho x_0)$,
so $\psi$ cannot single out the $\sigma^k$-direction.  (This is precisely the
mechanism used in the proof of the negative clause~P2 below; we record it here
as the content of (K4).)  In particular no \emph{single} existential reach
from outside a gadget certifies the translation direction, uniformly across
gadgets.
\end{proof}

\begin{remark}[On the wording ``unbounded displacement'']\label{rem:K4-honest}
The earlier informal statement of (K4) spoke of the per-gadget displacement of
$\sigma$ ``growing with $n$''.  In the present construction $\sigma$ acts by
the \emph{same} translation (by $1$) on every gadget, so what (K4) actually
provides --- and all that Sub-lemma~G5 uses --- is the
\emph{orientation-indistinguishability} of $\sigma^k$ from $\sigma^{-k}$ on an
unanchored gadget, established in Proposition~\ref{prop:K4}.  No metric growth
of displacement is needed.  Should a variant requiring literally growing
displacement be desired, one may replace the within-gadget path of $G_n$ by a
path subdivided so that the generator's induced minimal $P$-step on $G_n$ has
length $n+1$; this changes none of Lemmas~\ref{lem:shifts-are-auts}--\ref{lem:translation}
(the spine still couples the gadgets and the anchor still kills the reflection),
so $\AUT(\A)\cong\Z$ is preserved.  We do not need this refinement and so do
not carry it.
\end{remark}

\subsection*{Conclusion of G2}

\begin{corollary}\label{cor:G2}
The structure $\A$ of Definition~\ref{def:A} is a computably represented,
infinite, countable first-order structure with $\AUT(\A)\cong\Z$.  Hence
$\AUT(\A)$ is countable and nontrivial, the hypotheses of Theorems~2.2
and~2.3 are met, and $\A$ simultaneously enjoys properties
(K1)--(K4).  In particular, the linking spine creates no unexpected
automorphisms (Theorem~\ref{thm:aut} is an exact characterization) and does
not rigidify $\A$ past the availability of the diagonal clause
(Propositions~\ref{prop:K3}--\ref{prop:K4}).
\end{corollary}

\begin{proof}
Combine Theorem~\ref{thm:aut} with
Propositions~\ref{prop:K1}--\ref{prop:K4}.
\end{proof}

\section*{Sub-lemma (G5): joint consistency of the positive and negative clauses}

\subsection*{Setup and the apparent paradox}

We work inside the construction of Sub-lemma~G2 (the structure $\A$ of
Definition~\ref{def:A}, with $\AUT(\A)\cong\Z$ generated by $\sigma$).
$\Sigma^{in}_1$-definability has the exact meaning of the problem statement:
$f\colon\A\to\A$ is \emph{$\Sigma^{in}_1$-definable from a finite tuple
$\bar c$} if there is a family of existential formulas $(\phi_i)_{i\in\N}$
with, for all $x,y\in\A$,
\[
f(x)=y\quad\Longleftrightarrow\quad \exists i\;\bigl(\A\models\phi_i(\bar c,x,y)\bigr).
\tag{$\ast$}
\]

Two clauses are in play.
\begin{itemize}
\item[(P1)] \emph{Positive clause.} For \emph{every} $\pi\in\AUT(\A)$ there is
  \emph{some} finite tuple $\bar c_\pi$ such that $\pi$ is
  $\Sigma^{in}_1$-definable from $\bar c_\pi$.
\item[(P2)] \emph{Negative clause.} For \emph{every} finite tuple $\bar a$
  there is $\pi\in\AUT(\A)$ that is \emph{not} $\Sigma^{in}_1$-definable from
  $\bar a\cup\pi(\bar a)$.
\end{itemize}

The purpose of this sub-lemma is to determine whether (P1) and (P2) can hold
\emph{jointly} of $\A$, and to give the correct criterion for the negative
clause.  \textbf{The central outcome of this analysis is negative for the
present $\A$:} while (P1) holds, the structure $\A$ of Definition~\ref{def:A}
\emph{fails} (P2) outright.  We prove this, isolate exactly why, and state the
correct reduction that any successful construction must satisfy.  This corrects
an earlier draft of this sub-lemma, whose proof of (P2) relied on a false
``witness-locality'' claim.

\subsection*{(P1) holds: a defining tuple exists for each $\pi$}

\begin{lemma}[Positive clause]\label{lem:P1}
For every $\pi=\sigma^k\in\AUT(\A)$ there is a finite tuple $\bar c_\pi$ such
that $\pi$ is $\Sigma^{in}_1$-definable from $\bar c_\pi$.  Indeed
$\bar c_\pi$ may be taken to be the single point $(0,0)\in G_0$ (independently
of $k$), or even the empty tuple.
\end{lemma}

\begin{proof}
If $k=0$, $\pi=\mathrm{id}$ is defined by ``$x=y$''.  Assume $k\neq0$.
Take $\bar c:=\bigl((0,0)\bigr)$.  Along the spine, $(n,0)$ is the unique
point reached from $(0,0)$ by a directed $R$-path of length $n$, so
\[
\mathrm{Base}_n(\bar c,u):=\exists w_0\cdots w_n\Bigl(w_0=(0,0)\wedge w_n=u\wedge\bigwedge_{t=0}^{n-1}R(w_t,w_{t+1})\Bigr)
\]
defines $u=(n,0)$.  With $\theta_n^{(k)}$ from Proposition~\ref{prop:K2}, put
$\phi_n(\bar c,x,y):=\exists u\,(\mathrm{Base}_n(\bar c,u)\wedge\theta_n^{(k)}(u,x,y))$.
Then $\pi(x)=y\iff\exists n\,(\A\models\phi_n(\bar c,x,y))$, which is $(\ast)$.
(One can even drop the parameter: see Lemma~\ref{lem:leak} below, which shows
the orientation is recoverable spine\,$+$\,$D$ with no parameter at all.)
\end{proof}

\subsection*{(P2) fails for the present $\A$: the orientation leak}

The negative clause asks for an automorphism whose graph is \emph{not}
captured by any countable family of existential formulas over the parameters
$\bar a\cup\pi(\bar a)$.  For $\A$ this is impossible, because the directed
anchor $D$ on $G_0$ together with the position-aligned directed spine $R$
renders the orientation of \emph{every} gadget existentially definable with
\emph{no parameters at all}.

\begin{lemma}[Orientation leak]\label{lem:leak}
In the structure $\A$ of Definition~\ref{def:A}, the generator $\sigma$ (and
hence every $\sigma^k$) is $\Sigma^{in}_1$-definable from the \emph{empty}
tuple.
\end{lemma}

\begin{proof}
Define existential formulas $(\phi_n)_{n\ge0}$ by $\phi_0(x,y):=D(x,y)$ and,
for $n\ge1$,
\[
\phi_n(x,y):=\exists u_0\cdots u_n\,\exists v_0\cdots v_n\Bigl(
D(u_0,v_0)\wedge\bigwedge_{t=0}^{n-1}\bigl(R(u_t,u_{t+1})\wedge R(v_t,v_{t+1})\bigr)
\wedge x=u_n\wedge y=v_n\Bigr).
\]
We claim $\sigma(x)=y\iff\exists n\,(\A\models\phi_n(x,y))$.  Fix $x=(n,a)$.
The directed spine is a disjoint union of rays $(0,j)\to(1,j)\to\cdots$, so the
only $u_0,\dots,u_n$ with $R(u_t,u_{t+1})$ and $u_n=x$ are
$u_t=(t,a)$, forcing $u_0=(0,a)$; likewise $v_0=(0,b)$ with $v_n=y=(n,b)$.
The clause $D(u_0,v_0)=D((0,a),(0,b))$ holds iff $b=a+1$, i.e.\ $y=(n,a+1)=
\sigma(x)$.  Conversely $\phi_n((n,a),(n,a+1))$ holds with these witnesses.
For $m\neq n$ the formula $\phi_n$ fails on $x\in G_m$ (the $R$-chain of length
$n$ from a source $(0,\cdot)$ ends in $G_n$, not $G_m$).  Hence $(\ast)$ holds
with empty parameters.  (This is confirmed by a direct finite-model
computation on the truncation $0\le n\le4$, $|i|\le5$, where
$\bigvee_n\phi_n$ coincides with the graph of $\sigma$ on the interior.)
\end{proof}

\begin{corollary}[(P2) fails for $\A$]\label{cor:P2-fails}
For every finite tuple $\bar a$ and every $\pi=\sigma^k\in\AUT(\A)$, $\pi$
\emph{is} $\Sigma^{in}_1$-definable from $\bar a\cup\pi(\bar a)$ (indeed from
$\varnothing$).  Hence $\A$ does not witness the negative clause of
Question~1.2.
\end{corollary}

\begin{proof}
Immediate from Lemma~\ref{lem:leak}: adding parameters cannot destroy
definability, since the empty-parameter family $(\phi_n)$ already satisfies
$(\ast)$ for $\sigma^k$ (apply Lemma~\ref{lem:leak} to $\sigma$ and compose, or
repeat the proof with $D$ replaced by the $k$-fold $D$-chain).
\end{proof}

\begin{remark}[Diagnosis: where the earlier ``proof'' of (P2) broke]\label{rem:diagnosis}
An earlier draft argued (P2) by reflecting a far gadget $G_{n_0}$ disjoint
from $\bar c=\bar a\cup\pi(\bar a)$ via the partial automorphism $\rho$ of
Proposition~\ref{prop:K3}, then transporting a witnessing instance
$\phi_i(\bar c,x_0,\sigma x_0)$ to $\phi_i(\bar c,\rho x_0,\sigma^{-1}\rho x_0)$.
The transport step is \emph{invalid}: $\rho$ is only an automorphism of the
\emph{induced substructure} on $\bar c\cup G_{n_0}$, whereas the existential
witnesses of $\phi_i$ range over \emph{all} of $\A$.  Concretely, the
witnesses of the leaking formula $\phi_n$ of Lemma~\ref{lem:leak} run down the
spine from $G_{n_0}$ to the anchored gadget $G_0$ --- they leave
$\bar c\cup G_{n_0}$ entirely --- so $\rho$ does not act on them and the
quantifier-free matrix need not be preserved.  Properties (K3)/(K4) only assert
the (correct) fact that $\rho$ preserves the relations \emph{internal} to
$\bar c\cup G_{n_0}$; they say nothing about formulas whose witnesses escape
through the spine, and it is precisely such formulas that define $\sigma$.
\end{remark}

\subsection*{The correct criterion for the negative clause}

The fix is conceptual: full automorphisms of $\A$ are too rigid
($\AUT(\A)\cong\Z$, with trivial pointwise stabilizers), so invariance under
$\AUT(\A)$ cannot defeat $\Sigma^{in}_1$-definability.  The right tool is the
preservation of existential formulas under \emph{embeddings}, which are far
more plentiful than automorphisms.

\begin{lemma}[Embedding criterion for non-definability]\label{lem:embcrit}
Let $\B$ be any structure, $\bar c$ a finite tuple, and $\pi\colon\B\to\B$.
Suppose there is a self-embedding $h\colon\B\hookrightarrow\B$ (an injective
homomorphism preserving and reflecting all relations) with $h(\bar c)=\bar c$
pointwise and a point $x_0$ such that
\[
h\bigl(\pi(x_0)\bigr)\;\neq\;\pi\bigl(h(x_0)\bigr).
\]
Then $\pi$ is \emph{not} $\Sigma^{in}_1$-definable from $\bar c$.
\end{lemma}

\begin{proof}
Suppose $(\phi_i)$ witnessed $(\ast)$ for $\pi$ from $\bar c$.  Since
$\pi(x_0)=:y_0$, some $\phi_i$ has $\B\models\phi_i(\bar c,x_0,y_0)$.
Existential formulas are preserved by embeddings, so
$\B\models\phi_i(h\bar c,h x_0,h y_0)=\phi_i(\bar c,h x_0,h y_0)$.  By $(\ast)$
this forces $\pi(h x_0)=h y_0=h(\pi x_0)$, contradicting the hypothesis.
\end{proof}

\begin{remark}[What a successful construction must provide]\label{rem:spec}
By Lemma~\ref{lem:embcrit}, to witness Question~1.2 one needs a computably
$\AUT$-countable structure $\B$ such that:
\begin{enumerate}[label=(\roman*)]
\item (P1, hence Theorem~2.3) every $\pi\in\AUT(\B)$ is
  $\Sigma^{in}_1$-definable from \emph{some} finite tuple; and
\item (P2) for every finite $\bar a$ there are $\pi\in\AUT(\B)$ and a
  self-embedding $h$ of $\B$ fixing $\bar a\cup\pi(\bar a)$ pointwise with
  $h\circ\pi\neq\pi\circ h$.
\end{enumerate}
The structure $\A$ of Definition~\ref{def:A} \emph{provably fails} (ii):
Lemma~\ref{lem:leak} shows $\sigma$ is even parameter-free definable, so by
Lemma~\ref{lem:embcrit} \emph{no} self-embedding $h$ fixing $\varnothing$ can
fail to commute with $\sigma$ --- equivalently, the orientation that $h$ would
have to reverse is pinned down, parameter-free, by the spine\,$+$\,$D$.  The
two requirements (i) and (ii) are in genuine tension: any mechanism that
propagates orientation from finitely many anchored gadgets to all gadgets
(needed for (i) with a uniform shift) is equally available to the adversary's
existential formulas (defeating (ii)).  Resolving this tension --- making the
orientation \emph{determined} (Theorem~2.2) yet not \emph{existentially
definable} from external parameters (Theorem~2.3) --- is exactly the content of
the Theorem~2.2/2.3 gap, and the present $\A$ does \emph{not} inhabit it.
\end{remark}

\subsection*{Conclusion of G5}

\begin{theorem}[Status of joint consistency]\label{thm:G5}
For the structure $\A$ of Definition~\ref{def:A}: (P1) holds
(Lemma~\ref{lem:P1}) but (P2) \emph{fails} (Corollary~\ref{cor:P2-fails}).
Consequently $\A$ does not resolve Question~1.2.  The abstract reconciliation
of (P1) and (P2) is governed by the embedding criterion
(Lemma~\ref{lem:embcrit}): they are jointly consistent for a structure $\B$
exactly when $\B$ satisfies (i) and (ii) of Remark~\ref{rem:spec}, which $\A$
does not.
\end{theorem}

\begin{proof}
Lemma~\ref{lem:P1}, Corollary~\ref{cor:P2-fails}, Lemma~\ref{lem:embcrit}, and
Remark~\ref{rem:spec}.
\end{proof}

\begin{remark}[Honest summary]
There is no logical contradiction between (P1) and (P2) \emph{in general}:
(P1) speaks of a $\pi$-dependent tuple and definability via embeddings being
\emph{blockable}, while (P2) asks for an embedding fixing an
\emph{externally chosen} tuple that fails to commute with $\pi$
(Lemma~\ref{lem:embcrit}).  These quantify over different resources and are
not formally incompatible.  However, the \emph{specific} construction $\A$ of
Definition~\ref{def:A} does not realize this separation: its spine plus
directed anchor make orientation parameter-free definable, so (P2) is false
for it.  A correct counterexample to the effectivization of Theorem~2.2 must
replace the orientation-transmitting spine of Definition~\ref{def:A} by a
coupling that keeps $\AUT$ countable without rendering gadget orientation
existentially definable from external parameters; designing and verifying such
a coupling is left as the outstanding task (see the issue tracker).
\end{remark}

\end{document}
